\section{The interpolation operator and its normalization}\label{sec:preliminaries}
Put $Q=[-1,1]^2$ and let $\Pi_n$ be the real polynomials of total degree
at most $n$. Here $\delta_{AB}$ is one if $A=B$ and zero otherwise.
The cardinal polynomials satisfy $L_B(A)=\delta_{AB}$ for
$A,B\in\mathcal P_n$, and
\[
 I_nf(z)=\sum_{B\in\mathcal P_n} f(B)L_B(z).
\]
The triangle inequality gives $\|I_n\|\leq\max_Q\lambda_n$; equality
follows by prescribing at the nodes the signs of the cardinal values
at a maximizing point and extending these finitely many values
continuously to $Q$. This proves the operator-norm interpretation used
in \eqref{eq:lebesgue-inequality}.

We use the Chebyshev polynomials $T_r(\cos t)=\cos(rt)$. Write
$\theta_j=j\pi/n$, $\phi_k=k\pi/(n+1)$ and
\[
 \varepsilon_j=\begin{cases}1/2,&j=0,n,\\1,&\text{otherwise},\end{cases}
 \qquad
 \eta_k=\begin{cases}1/2,&k=0,n+1,\\1,&\text{otherwise},\end{cases}
 \qquad w_{jk}=\frac{2\varepsilon_j\eta_k}{n(n+1)}.
\]
For the kernel formula below, set $\nu_0=1$ and $\nu_r=2$ for $r>0$;
these are the mode multiplicities. Define
\[
 K_n(z,B)=\sum_{p+q\leq n}\nu_p\nu_q
 T_p(z_1)T_p(B_1)T_q(z_2)T_q(B_2).
\]

\begin{lemma}[Cardinal kernel]\label{lem:cardinal-kernel}
The set $\mathcal P_n$ is unisolvent for $\Pi_n$, and for
$B=(\cos\theta_j,\cos\phi_k)\in\mathcal P_n$,
\begin{equation}\label{eq:cardinal-kernel}
 L_B(z)=w_{jk}\bigl[K_n(z,B)-T_n(z_1)T_n(B_1)\bigr].
\end{equation}
In particular the coefficient of the highest first-coordinate mode is
one.
\end{lemma}
\begin{proof}
For $M\geq1$ let
\[
 Q_M(g)=\frac1M\sum_{j=0}^M{}'g(j\pi/M),
\]
where the prime gives weight $1/2$ to each endpoint. Finite geometric
sums show that $Q_M(\cos rt)=0$ for $1\leq r\leq2M-1$ and
$Q_M(1)=Q_M(\cos2Mt)=1$. Consequently
\[
 Q_M(\cos pt\cos rt)=\delta_{pr}h_p^{(M)},\qquad
 h_p^{(M)}=\begin{cases}1,&p=0,M,\\1/2,&0<p<M.\end{cases}
\]
Multiplication by $(-1)^j$ sends the mode $p$ to $M-p$ on this grid.
In the checkerboard scalar product the identity
$2\mathbf1_{j+k\,\mathrm{even}}=1+(-1)^{j+k}$ splits the Gram matrix
of the modes $T_p(x)T_q(y)$, $p+q\leq n$, into a tensor product part
and an alternating part. The alternating part could be nonzero only
if $p+r=n$ and $q+s=n+1$. Their sum contradicts
$(p+q)+(r+s)\leq2n$, so that part vanishes.
The Gram matrix is therefore diagonal and positive, with diagonal
$h_p^{(n)}h_q^{(n+1)}$.

There are $(n+1)(n+2)/2$ checkerboard points, the same number as the
modes. The square evaluation matrix is invertible, proving unisolvence.
The inverse Gram coefficients are $\nu_p\nu_q$ except at $(p,q)=(n,0)$,
where the coefficient is one. The discrete reproducing kernel, multiplied
by $w_{jk}$, therefore has the cardinal values and gives
\eqref{eq:cardinal-kernel}.
\end{proof}

The same correction appears in the angular kernel
\begin{equation}\label{eq:angular-kernel}
 G_n(u,v)=\sum_{p+q\leq n}\nu_p\nu_q\cos(pu)\cos(qv)-\cos(nu).
\end{equation}
At a corner it gives the exact signed row sums; in the interior its
reflections give the cancellation comparison. Thus the two arguments
start from one normalized kernel.



The two nonnode corners in Theorem~\ref{thm:main-location} have equal
Lebesgue values by symmetry. For even $n$, reflection in the first
coordinate sends $j$ to $n-j$ and preserves the checkerboard. For odd
$n$, reflection in the second coordinate sends $k$ to $n+1-k$ and
preserves it. Uniqueness of interpolation carries these symmetries to
the cardinal polynomials and hence to $\lambda_n$.

\begin{remark}
All formulas involving trigonometric quotients are interpreted by their
continuous polynomial limits at removable singularities. The maximization
and certificate domains include all nodes, edges and collisions of angular
coordinates.
\end{remark}
